\section{Red de sensores}

La red de sensado está distribuida entre los tres ATmega328P. El Maestro solicita las mediciones de los sensores conectados a los Slaves mediante I2C y, además, accede directamente al sensor ToF VL53L0X. Esta distribución permite que cada microcontrolador atienda localmente sus dispositivos y que el Maestro concentre la toma de decisiones del proceso.

\subsection{Sensor infrarrojo de entrada}
El sensor infrarrojo está conectado al Slave 2 en PC0/A0. Su salida es activa en bajo: un nivel lógico bajo representa detección del vehículo y un nivel alto representa espacio libre. El Slave 2 actualiza la lectura aproximadamente cada 10 ms y mantiene un contador de muestras para que el Maestro pueda distinguir datos nuevos.

En estado de reposo, el Maestro consulta el sensor cada 100 ms. Para reducir activaciones causadas por cambios aislados, la lógica automática requiere cinco detecciones consecutivas antes de iniciar el ingreso del vehículo. Una vez confirmada la presencia, el Maestro ordena abrir la puerta de entrada.

\subsection{HC-SR04 y nivel del depósito}
El HC-SR04 está conectado al Slave 1, con ECHO en PC2/A2 y TRIG en PC3/A3. El Slave 1 realiza la adquisición local y entrega al Maestro el ancho del pulso de eco. El Maestro convierte posteriormente dicho tiempo a distancia.

Durante la etapa asociada al suministro de líquido, la primera medición válida se almacena como distancia de referencia. Después se abre el servo de agua a 90 grados y se continúa monitoreando el depósito. Cuando la distancia medida aumenta 20 mm respecto de la referencia, el sistema cierra el servo y reanuda el desplazamiento del automóvil. Por tanto, la medición ultrasónica no se utiliza únicamente para visualización, sino también como condición de control del proceso.

\subsection{VL53L0X}
El VL53L0X, identificado en la arquitectura como Slave 3, está conectado directamente al bus I2C del Maestro y utiliza la dirección 0x29. La línea XSHUT está conectada a PB2/D10 y permite mantener el sensor apagado o reiniciarlo físicamente. En la implementación actual, el Maestro inicializa el sensor cuando comienza el transporte del vehículo y realiza mediciones de distancia durante distintas etapas.

El código utiliza el VL53L0X como referencia de posición a lo largo del recorrido. La zona de suministro se identifica entre 265 y 270 mm, la estación de cepillado entre 125 y 130 mm, y la condición de finalización se detecta cuando la distancia es menor o igual a 70 mm. Las mediciones consideradas físicamente válidas por la lógica se encuentran entre 20 y 2000 mm. Estos umbrales corresponden a la implementación actual y deberán contrastarse con las pruebas finales de la maqueta.
